% GENERATED FILE. Edit canonical lessons, then run npm run generate:books.
% canonical-source-hash: fnv1a64:a1c0cb3c7884238f
% canonical-lessons: SA-C523-natah, SA-C523-abhinayati, SA-C523-vidusakah, SA-C523-nepathyam, SA-C523-abhinayah

\chapter{A Play in the Village}
\label{ch:sa-a-play-in-the-village}
\hlchaptermodality{523}

\begin{chapteropening}
\textbf{By the end of this chapter:} \emph{I can talk about a play and the people in it.}

Complete the last lesson of chapter 523: I can talk about a play and the people in it.
\end{chapteropening}

\section[\texorpdfstring{\sk{नटः} (naṭaḥ)}{नटः (naṭaḥ)}]{\sk{नटः} (naṭaḥ) — an actor}

\label{lesson:SA-C523-natah}



\begin{quote}
Before the new one: say the Sanskrit for imitates, then the Sanskrit for clapping.
\end{quote}

\subsection*{You'll want to know: \sk{नटः}}
\textbf{\sk{नटः}} — \emph{naṭaḥ} — \textquotedblleft{}an actor\textquotedblright{}.

\begin{quote}
\sk{मम} \sk{मित्रं} \sk{नाटके} \sk{नटः} \sk{अस्ति।}
\end{quote}

\begin{quote}
\emph{mama mitraṃ nāṭake naṭaḥ asti.}
\end{quote}

\begin{quote}
My friend is an actor in the play.
\end{quote}

A masculine noun.

\begin{cousinweb}
\textbf{\sk{नटः}} (\emph{naṭaḥ}) comes from the root √naṭ, \textquotedblleft{}to dance\textquotedblright{}, as MW derives it; MW gives \emph{naṭa} as an actor, a dancer, a mime (MW).
\end{cousinweb}

\begin{grammarlens}[title={What you've built}]
One more word for describing what you lived and felt.
\end{grammarlens}

\subsection*{Your turn}
\begin{itemize}
\raggedright
  \item \emph{Say it:} \emph{naṭaḥ}
  \item \emph{Say it:} \emph{naṭaḥ}, once more
  \item \emph{Say it:} \emph{mama mitraṃ nāṭake naṭaḥ asti.} — My friend is an actor in the play.
  \item \emph{Recall:} say the Sanskrit for imitates, then the Sanskrit for clapping, then say \emph{naṭaḥ} again
\end{itemize}

\begin{tcolorbox}[breakable,colback=teal!4,colframe=teal!35!black,title={Before you move on}]
What does \sk{नटः} mean? (\textquotedblleft{}An actor\textquotedblright{}.) Say it once more.
\end{tcolorbox}
\section[\texorpdfstring{\sk{अभिनयति} (abhinayati)}{अभिनयति (abhinayati)}]{\sk{अभिनयति} (abhinayati) — acts a part}

\label{lesson:SA-C523-abhinayati}



\begin{quote}
Before the new one: say the Sanskrit for clapping, then the Sanskrit for an actor.
\end{quote}

\subsection*{You'll want to know: \sk{अभिनयति}}
\textbf{\sk{अभिनयति}} — \emph{abhinayati} — \textquotedblleft{}acts a part\textquotedblright{}.

\begin{quote}
\sk{सः} \sk{नाटके} \sk{चौरम्} \sk{अभिनयति।}
\end{quote}

\begin{quote}
\emph{saḥ nāṭake cauram abhinayati.}
\end{quote}

\begin{quote}
He plays a thief in the play.
\end{quote}

Past (he …): \sk{सः} \sk{चौरम्} \sk{अभिनीतवान्}, \emph{saḥ cauram abhinītavān}.

\begin{cousinweb}
\textbf{\sk{अभिनयति}} (\emph{abhinayati}) is built on \emph{abhi}, \textquotedblleft{}towards\textquotedblright{}, and the root √nī, \textquotedblleft{}to lead\textquotedblright{}, which also gives \emph{nayati}; MW gives \emph{abhinī} as to bring near, and to represent dramatically, to act (MW).
\end{cousinweb}

\begin{grammarlens}[title={What you've built}]
One more word for describing what you lived and felt.
\end{grammarlens}

\subsection*{Your turn}
\begin{itemize}
\raggedright
  \item \emph{Say it:} \emph{abhinayati}
  \item \emph{Say it:} \emph{abhinayati}, once more
  \item \emph{Say it:} \emph{saḥ nāṭake cauram abhinayati.} — He plays a thief in the play.
  \item \emph{Recall:} say the Sanskrit for clapping, then the Sanskrit for an actor, then say \emph{abhinayati} again
\end{itemize}

\begin{tcolorbox}[breakable,colback=teal!4,colframe=teal!35!black,title={Before you move on}]
What does \sk{अभिनयति} mean? (\textquotedblleft{}Acts a part\textquotedblright{}.) Say it once more.
\end{tcolorbox}
\section[\texorpdfstring{\sk{विदूषकः} (vidūṣakaḥ)}{विदूषकः (vidūṣakaḥ)}]{\sk{विदूषकः} (vidūṣakaḥ) — a jester}

\label{lesson:SA-C523-vidusakah}



\begin{quote}
Before the new one: say the Sanskrit for an actor, then the Sanskrit for acts a part.
\end{quote}

\subsection*{You'll want to know: \sk{विदूषकः}}
\textbf{\sk{विदूषकः}} — \emph{vidūṣakaḥ} — \textquotedblleft{}a jester\textquotedblright{}.

\begin{quote}
\sk{विदूषकं} \sk{दृष्ट्वा} \sk{सर्वे} \sk{हसन्ति।}
\end{quote}

\begin{quote}
\emph{vidūṣakaṃ dṛṣṭvā sarve hasanti.}
\end{quote}

\begin{quote}
Everyone laughs at the sight of the jester.
\end{quote}

A masculine noun.

\begin{cousinweb}
\textbf{\sk{विदूषकः}} (\emph{vidūṣakaḥ}) is built on \emph{vi} and the causative of the root √duṣ, \textquotedblleft{}to spoil, to defile\textquotedblright{}, with the agent ending \emph{-aka}; MW gives \emph{vidūṣaka} first as defiling, then as a jester, above all the witty friend of the hero in a play (MW).
\end{cousinweb}

\begin{grammarlens}[title={What you've built}]
One more word for describing what you lived and felt.
\end{grammarlens}

\subsection*{Your turn}
\begin{itemize}
\raggedright
  \item \emph{Say it:} \emph{vidūṣakaḥ}
  \item \emph{Say it:} \emph{vidūṣakaḥ}, once more
  \item \emph{Say it:} \emph{vidūṣakaṃ dṛṣṭvā sarve hasanti.} — Everyone laughs at the sight of the jester.
  \item \emph{Recall:} say the Sanskrit for an actor, then the Sanskrit for acts a part, then say \emph{vidūṣakaḥ} again
\end{itemize}

\begin{tcolorbox}[breakable,colback=teal!4,colframe=teal!35!black,title={Before you move on}]
What does \sk{विदूषकः} mean? (\textquotedblleft{}A jester\textquotedblright{}.) Say it once more.
\end{tcolorbox}
\section[\texorpdfstring{\sk{नेपथ्यम्} (nepathyam)}{नेपथ्यम् (nepathyam)}]{\sk{नेपथ्यम्} (nepathyam) — backstage}

\label{lesson:SA-C523-nepathyam}



\begin{quote}
Before the new one: say the Sanskrit for acts a part, then the Sanskrit for a jester.
\end{quote}

\subsection*{You'll want to know: \sk{नेपथ्यम्}}
\textbf{\sk{नेपथ्यम्}} — \emph{nepathyam} — \textquotedblleft{}backstage\textquotedblright{}.

\begin{quote}
\sk{नटः} \sk{नेपथ्ये} \sk{वस्त्रं} \sk{परिवर्तयति।}
\end{quote}

\begin{quote}
\emph{naṭaḥ nepathye vastraṃ parivartayati.}
\end{quote}

\begin{quote}
The actor changes his clothes backstage.
\end{quote}

A neuter noun.

\begin{cousinweb}
\textbf{\sk{नेपथ्यम्}} (\emph{nepathyam}) is the word MW gives for an ornament, a costume, and in drama the place behind the stage, the tiring-room; MW gives no derivation, so its origin is uncertain (MW).
\end{cousinweb}

\begin{grammarlens}[title={What you've built}]
One more word for describing what you lived and felt.
\end{grammarlens}

\subsection*{Your turn}
\begin{itemize}
\raggedright
  \item \emph{Say it:} \emph{nepathyam}
  \item \emph{Say it:} \emph{nepathyam}, once more
  \item \emph{Say it:} \emph{naṭaḥ nepathye vastraṃ parivartayati.} — The actor changes his clothes backstage.
  \item \emph{Recall:} say the Sanskrit for acts a part, then the Sanskrit for a jester, then say \emph{nepathyam} again
\end{itemize}

\begin{tcolorbox}[breakable,colback=teal!4,colframe=teal!35!black,title={Before you move on}]
What does \sk{नेपथ्यम्} mean? (\textquotedblleft{}Backstage\textquotedblright{}.) Say it once more.
\end{tcolorbox}
\section[\texorpdfstring{\sk{अभिनयः} (abhinayaḥ)}{अभिनयः (abhinayaḥ)}]{\sk{अभिनयः} (abhinayaḥ) — acting}

\label{lesson:SA-C523-abhinayah}



\begin{quote}
Before the new one: say the Sanskrit for a jester, then the Sanskrit for backstage.
\end{quote}

\subsection*{You'll want to know: \sk{अभिनयः}}
\textbf{\sk{अभिनयः}} — \emph{abhinayaḥ} — \textquotedblleft{}acting\textquotedblright{}.

\begin{quote}
\sk{तस्य} \sk{अभिनयः} \sk{मह्यं} \sk{रोचते।}
\end{quote}

\begin{quote}
\emph{tasya abhinayaḥ mahyaṃ rocate.}
\end{quote}

\begin{quote}
I like his acting.
\end{quote}

A masculine noun.

\begin{cousinweb}
\textbf{\sk{अभिनयः}} (\emph{abhinayaḥ}) is built on \emph{abhi} and the root √nī, \textquotedblleft{}to lead\textquotedblright{}; MW gives \emph{abhinaya} as acting, dramatic action, the showing of a feeling by look and gesture (MW).
\end{cousinweb}

\begin{grammarlens}[title={What you've built}]
One more word for describing what you lived and felt.
\end{grammarlens}

\subsection*{Your turn}
\begin{itemize}
\raggedright
  \item \emph{Say it:} \emph{abhinayaḥ}
  \item \emph{Say it:} \emph{abhinayaḥ}, once more
  \item \emph{Say it:} \emph{tasya abhinayaḥ mahyaṃ rocate.} — I like his acting.
  \item \emph{Recall:} say the Sanskrit for a jester, then the Sanskrit for backstage, then say \emph{abhinayaḥ} again
\end{itemize}

\begin{tcolorbox}[breakable,colback=teal!4,colframe=teal!35!black,title={Before you move on}]
What does \sk{अभिनयः} mean? (\textquotedblleft{}Acting\textquotedblright{}.) Say it once more.
\end{tcolorbox}
